\section{Discussion and limitations}
\label{sec:discussion}

\paragraph{Decoder-free optical design.} Our results show that improving depth discrimination at the measurement stage can benefit different reconstruction methods without coupling optical training to a decoder. The learned plate reduces depth error under TV and two PnP decoders, despite being optimised without specimen examples or reconstruction feedback. These gains are not tied to a single decoder, although the two PnP methods share an underlying framework. Unlike end-to-end optical design, which optimises a particular optics--decoder combination, our approach directly targets depth discrimination in the measurements. We do not establish superiority over end-to-end learning or transfer to arbitrary reconstruction methods.

\paragraph{Depth recovery and measurement ambiguity.} The learned plate improves depth placement, and conformal calibration turns the improvement into narrower intervals with an expected-risk guarantee under exchangeability. It does not make neuronal topology identifiable: at eight-fold compression no plate determines dense neuronal windows from the measurements alone (\cref{sec:res:guarantee}). In supplementary experiments, conformal connectivity calls on these reconstructions missed their nominal error rate, and query-adaptive direction selection did not meet its prespecified criterion. Connections drawn from such reconstructions remain the prior's choice.

\paragraph{Physical assumptions and implementation.} Our results are based entirely on simulations of linear fluorescence in a homogeneous, non-scattering medium with idealised optical components. Scattering, aberrations, and other experimental imperfections may alter the depth-dependent illumination patterns and the depth discrimination they provide. Although the alignment simulations show that the learned gains can survive moderate perturbations when the optical model is updated, uncalibrated misalignment can eliminate those gains. Experimental validation therefore requires accurate characterisation of the realised optical system and its measurement operator.

\paragraph{Data and statistical limitations.} Our evaluation uses 60 neuronal crops from six cells across three animals. Crops from the same cell are correlated, and the cell-bootstrap confidence intervals remain approximate with only six cells. The objective comparison and depth-interval evaluation use 12 prespecified crops, while optical training was repeated with two seeds. The denoisers were trained on a fourth animal, with conformal calibration performed on a held-out stack from that animal. Since the test volumes come from different animals, exchangeability is not guaranteed, and empirical interval coverage does not establish a cross-animal guarantee. Broader evaluation across specimens and imaging conditions is needed to assess generalisation.
